\documentclass[11pt]{article}
\usepackage[margin=1.1in]{geometry}
\usepackage{amsmath,amssymb,amsthm,mathtools,booktabs,array,tabularx,longtable}
\usepackage[dvipsnames]{xcolor}
\usepackage[colorlinks=true,linkcolor=MidnightBlue,citecolor=MidnightBlue,urlcolor=MidnightBlue]{hyperref}
\usepackage{fancyhdr}
\pagestyle{fancy}\fancyhf{}\renewcommand{\headrulewidth}{0pt}
\fancyfoot[R]{\thepage}
\fancyfoot[L]{\scriptsize\color{gray}Generated by the Transport Lab from approved parts of its store; not a deposited paper. Papers of record: \href{https://zenodo.org/communities/transport-papers}{Zenodo, Transport Papers}}
\fancypagestyle{plain}{\fancyhf{}\fancyfoot[R]{\thepage}\fancyfoot[L]{\scriptsize\color{gray}Generated by the Transport Lab from approved parts of its store; not a deposited paper. Papers of record: \href{https://zenodo.org/communities/transport-papers}{Zenodo, Transport Papers}}\renewcommand{\headrulewidth}{0pt}}
\newtheorem{theorem}{Theorem}[section]
\newtheorem{proposition}[theorem]{Proposition}
\newtheorem{lemma}[theorem]{Lemma}
\newtheorem{corollary}[theorem]{Corollary}
\newtheorem{conjecture}[theorem]{Conjecture}
\newtheorem{refutation}[theorem]{Refutation}
\newtheorem{prediction}[theorem]{Prediction}
\newtheorem{principle}[theorem]{Principle}
\newtheorem{observation}[theorem]{Observation}
\newtheorem{axiom}[theorem]{Axiom}
\theoremstyle{definition}
\newtheorem{definition}[theorem]{Definition}
\newtheorem{fact}[theorem]{Fact}
\newtheorem{assembled}[theorem]{Assembly}
\newtheorem{postulate}[theorem]{Postulate}
\newtheorem{premise}[theorem]{Premise}
\newtheorem{hypothesis}[theorem]{Hypothesis}
\newtheorem{remark}[theorem]{Remark}
\newtheorem{example}[theorem]{Example}
\theoremstyle{plain}
\setlength{\marginparwidth}{0.8in}
\newcommand{\changedmark}[2]{\marginpar{\raggedright\scriptsize\color{#1}#2}}
\providecommand{\Aa}{}\renewcommand{\Aa}{\mathcal{A}}
\providecommand{\Ad}{}\renewcommand{\Ad}{\widetilde{\mathrm{Ad}}}
\providecommand{\Aev}{}\renewcommand{\Aev}{\mathcal{A}^{+}}
\providecommand{\Aod}{}\renewcommand{\Aod}{\mathcal{A}^{-}}
\providecommand{\B}{}\renewcommand{\B}{\mathrm{B}}
\providecommand{\C}{}\renewcommand{\C}{\mathbb{C}}
\providecommand{\CAR}{}\renewcommand{\CAR}{\mathrm{CAR}}
\providecommand{\CS}{}\renewcommand{\CS}{{\mathbb{S}^{+}}}
\providecommand{\CSm}{}\renewcommand{\CSm}{{\mathbb{S}^{-}}}
\providecommand{\Cl}{}\renewcommand{\Cl}{\mathrm{Cl}}
\providecommand{\Con}{}\renewcommand{\Con}{\mathcal{C}}
\providecommand{\Dl}{}\renewcommand{\Dl}{\Delta}
\providecommand{\EN}{}\renewcommand{\EN}{E_{N}}
\providecommand{\End}{}\renewcommand{\End}{\mathrm{End}}
\providecommand{\F}{}\renewcommand{\F}{\mathbb{F}}
\providecommand{\G}{}\renewcommand{\G}{\Gamma}
\providecommand{\HH}{}\renewcommand{\HH}{\mathbb{H}}
\providecommand{\IS}{}\renewcommand{\IS}{I\sigma_3}
\providecommand{\J}{}\renewcommand{\J}{\mathbf{J}}
\providecommand{\LH}{}\renewcommand{\LH}{L_{H}}
\providecommand{\Mat}{}\renewcommand{\Mat}{\mathrm{M}}
\providecommand{\Mo}{}\renewcommand{\Mo}{M^{\circ}}
\providecommand{\Mp}{}\renewcommand{\Mp}{M^{+}}
\providecommand{\Msun}{}\renewcommand{\Msun}{M_{\odot}}
\providecommand{\Mt}{}\renewcommand{\Mt}{\widetilde{M}}
\providecommand{\Ncone}{}\renewcommand{\Ncone}{\mathcal{N}}
\providecommand{\Neff}{}\renewcommand{\Neff}{N_{\mathrm{eff}}}
\providecommand{\PCS}{}\renewcommand{\PCS}{\mathbb{P}{\mathbb{S}^{+}}}
\providecommand{\PT}{}\renewcommand{\PT}{\mathbb{PT}}
\providecommand{\Pin}{}\renewcommand{\Pin}{\mathrm{Pin}}
\providecommand{\Q}{}\renewcommand{\Q}{\mathbb{Q}}
\providecommand{\R}{}\renewcommand{\R}{\mathbb{R}}
\providecommand{\Ric}{}\renewcommand{\Ric}{\mathrm{Ric}}
\providecommand{\SBH}{}\renewcommand{\SBH}{S_{\mathrm{BH}}}
\providecommand{\Sc}{}\renewcommand{\Sc}{\mathbb{S}}
\providecommand{\Sch}{}\renewcommand{\Sch}{\mathrm{P}}
\providecommand{\Sig}{}\renewcommand{\Sig}{\Sigma}
\providecommand{\Som}{}\renewcommand{\Som}{S_\omega}
\providecommand{\Sp}{}\renewcommand{\Sp}{\mathbb{S}}
\providecommand{\Sph}{}\renewcommand{\Sph}{\mathrm{S}}
\providecommand{\Spi}{}\renewcommand{\Spi}{S_\pi}
\providecommand{\Spin}{}\renewcommand{\Spin}{\mathrm{Spin}}
\providecommand{\Sth}{}\renewcommand{\Sth}{S_{\mathrm{th}}}
\providecommand{\T}{}\renewcommand{\T}{\mathbb{T}}
\providecommand{\TH}{}\renewcommand{\TH}{T_{\mathrm{H}}}
\providecommand{\Tr}{}\renewcommand{\Tr}{\operatorname{Tr}}
\providecommand{\UHF}{}\renewcommand{\UHF}{\mathrm{UHF}}
\providecommand{\Z}{}\renewcommand{\Z}{\mathbb{Z}}
\providecommand{\Zc}{}\renewcommand{\Zc}{\mathrm{Z}}
\providecommand{\cmg}{}\renewcommand{\cmg}{\mathrm{cm^{2}\,g^{-1}}}
\providecommand{\dS}{}\renewcommand{\dS}{\mathrm{dS}}
\providecommand{\dsc}{}\renewcommand{\dsc}{\mathrm{dsc}}
\providecommand{\dtPG}{}\renewcommand{\dtPG}{dt}
\providecommand{\dtS}{}\renewcommand{\dtS}{dt_{\mathrm{S}}}
\providecommand{\ellv}{}\renewcommand{\ellv}{\ell}
\providecommand{\er}{}\renewcommand{\er}{e_r}
\providecommand{\et}{}\renewcommand{\et}{e_t}
\providecommand{\ev}{}\renewcommand{\ev}{\mathrm{ev}}
\providecommand{\gl}{}\renewcommand{\gl}{\mathfrak{gl}}
\providecommand{\hatV}{}\renewcommand{\hatV}{\widehat{V}}
\providecommand{\hb}{}\renewcommand{\hb}{\bar{h}}
\providecommand{\hf}{}\renewcommand{\hf}{\hat{f}}
\providecommand{\id}{}\renewcommand{\id}{\mathrm{id}}
\providecommand{\kpc}{}\renewcommand{\kpc}{\mathrm{kpc}}
\providecommand{\lP}{}\renewcommand{\lP}{\ell_{P}}
\providecommand{\mD}{}\renewcommand{\mD}{\mathbf{m}_{D}}
\providecommand{\mb}{}\renewcommand{\mb}{\bar m}
\providecommand{\mnu}{}\renewcommand{\mnu}{\Sigma m_{\nu}}
\providecommand{\nInf}{}\renewcommand{\nInf}{n_\infty}
\providecommand{\nO}{}\renewcommand{\nO}{n_0}
\providecommand{\ninf}{}\renewcommand{\ninf}{n_{\infty}}
\providecommand{\nz}{}\renewcommand{\nz}{n_{0}}
\providecommand{\proj}{}\renewcommand{\proj}{\operatorname{proj}}
\providecommand{\pv}{}\renewcommand{\pv}{\pi}
\providecommand{\rad}{}\renewcommand{\rad}{\mathrm{rad}}
\providecommand{\sgn}{}\renewcommand{\sgn}{\operatorname{sgn}}
\providecommand{\shA}{}\renewcommand{\shA}{\mathcal{A}}
\providecommand{\shC}{}\renewcommand{\shC}{\mathcal{C}}
\providecommand{\shE}{}\renewcommand{\shE}{\mathcal{E}}
\providecommand{\shJ}{}\renewcommand{\shJ}{\mathcal{J}}
\providecommand{\shK}{}\renewcommand{\shK}{\mathcal{K}}
\providecommand{\shL}{}\renewcommand{\shL}{L}
\providecommand{\shW}{}\renewcommand{\shW}{\mathcal{W}}
\providecommand{\sll}{}\renewcommand{\sll}{\mathfrak{sl}}
\providecommand{\so}{}\renewcommand{\so}{\mathfrak{so}}
\providecommand{\su}{}\renewcommand{\su}{\mathfrak{su}}
\providecommand{\supp}{}\renewcommand{\supp}{\operatorname{supp}}
\providecommand{\tr}{}\renewcommand{\tr}{\operatorname{tr}}
\providecommand{\uu}{}\renewcommand{\uu}{\mathfrak{u}}
\providecommand{\br}{}\renewcommand{\br}[1]{\langle #1\rangle}
\providecommand{\bra}{}\renewcommand{\bra}[1]{\langle #1|}
\providecommand{\file}{}\renewcommand{\file}[1]{\texttt{#1}}
\providecommand{\hp}{}\renewcommand{\hp}[2]{( #1 \,|\, #2 )}
\providecommand{\ip}{}\renewcommand{\ip}[2]{\langle #1, #2\rangle}
\providecommand{\ket}{}\renewcommand{\ket}[1]{|#1\rangle}
\providecommand{\rel}{}\renewcommand{\rel}[2]{D\bigl(#1 \,\|\, #2\bigr)}
\providecommand{\rev}{}\renewcommand{\rev}[1]{\widetilde{#1}}
\providecommand{\slashed}{}\renewcommand{\slashed}[1]{{#1}\mkern-10mu/}
\title{\bfseries The algebra core}
\author{Leon Butler\\ \normalsize\href{mailto:lbutler.math@gmail.com}{\texttt{lbutler.math@gmail.com}}\\ \normalsize ORCID: \href{https://orcid.org/0009-0006-4449-5617}{0009-0006-4449-5617}}
\date{Generated by the Transport Lab, 2 November 2026}
\begin{document}
\maketitle
\begin{abstract}
\normalsize
Every result in the Transport Papers is stated in one algebra: the geometric algebra of a space of directions with a rule for their square lengths. This trail walks through its core, from that one rule to rotors, the classification of the real algebras, and the fermion Fock space.

Products of basis directions form a basis of the whole algebra, even when some directions have zero length. Sandwiching between directions carries out every reflection and rotation, and each rotation comes from exactly two rotors, $R$ and $-R$.

Pairs of light-like directions behave exactly like fermion creation and annihilation operators.
\end{abstract}
\newpage
\noindent \textbf{Authorship and disclosure. }The founding idea of the series is the author's: that division by zero in a Clifford algebra should be read not as a prohibition but as a transport into a higher-dimensional algebra, and that the transported division need not be accompanied by an inverse operation. The direction of investigation and the critical review of the results are likewise the author's. The mathematical development and the drafting of the text were carried out in extended collaboration with Claude Fable 5 (Anthropic), an AI system.

\noindent\fbox{\parbox{\dimexpr\linewidth-2\fboxsep-2\fboxrule}{\textbf{This is not a deposited paper.} It was generated by the Transport Lab on 2 November 2026 and assembled from approved parts of its store. It follows results across several papers of the series; each is stated with its proof and credited to the paper that first proved it. The official versions of those papers are on Zenodo, in the Transport Papers community, \url{https://zenodo.org/communities/transport-papers}.}}
\newpage

\section*{Introduction}
The walk starts with the defining rule, builds up grades and symmetries, and ends where the algebra meets quantum theory.

\section{Clifford algebras and quadratic forms}

These results build the geometric algebra from one rule and show what it contains.

The algebra has a single defining rule: a vector $v$ multiplied by itself gives its square length $Q(v)$, a plain number.

Applying that rule to a sum of two vectors $u$ and $v$ gives \[ uv + vu = 2B(u, v), \] where $B$ is the dot product that $Q$ defines. So the product $uv$ is the scalar $B(u, v)$ plus a part that changes sign when the order is swapped.

\begin{theorem}[The product of two vectors]\label{c:std-algebra.vector-squares}
Multiplying two arrows gives a number plus something new. The number is their ordinary dot product. The new part changes sign when the order is swapped, and turns out to be the oriented plane the two arrows span.
\emph{First proved in} the algebra core of the Transport Lab, proved there with an exact check; standard, see C. Doran and A. Lasenby, Geometric Algebra for Physicists, Cambridge University Press (2003), Ch. 4 (foundations of geometric algebra).
\end{theorem}
\noindent\emph{Rests on:} \ref{c:std-algebra.polarisation}.

Reordering costs only a sign and a scalar, and a repeated vector collapses to a number, so every product reduces to products $e_I$ of distinct basis vectors in increasing order.

In $n$ dimensions these blades, from $1$ up to the full volume, form a basis, so \[ \dim \Cl(V, Q) = 2\^{}{n} \] This holds when some directions have zero length, which matters because the series works in exactly such degenerate spaces.

\begin{theorem}[The blade basis, degenerate forms included]\label{c:std-algebra.blade-basis}
Every element of the algebra is a combination of products of distinct basis arrows, and no such product can be written in terms of the others. In $n$ dimensions there are $2^{n}$ of them: the number, the arrows, the planes, the volumes and so on. This stays true even when some directions have zero length.
\emph{First proved in} the algebra core of the Transport Lab, proved there with an exact check; standard, see C. Chevalley, The Algebraic Theory of Spinors, Columbia University Press (1954), Ch. II (the Clifford algebra of a quadratic space).
\end{theorem}
\noindent\emph{Rests on:} \ref{c:std-algebra.existence-uniqueness}, \ref{c:std-algebra.vector-squares}.

Each blade has a grade $k$, the number of vectors in it, so every element splits into scalars, vectors, bivectors and so on.

Three ways of flipping an element act on grade $k$ by a sign: the grade involution, reversion (writing every product backwards) and their composite. Reversion, written with a tilde, acts by \[ \tilde{x} = (-1)\^{}{k(k-1)/2} x \] on grade $k$, and it is the one that undoes rotations later in the walk.

\begin{theorem}[The three involutions on each grade]\label{c:std-algebra.involution-properties}
The three basic ways of "flipping" an element each multiply every grade by plus or minus one, according to a simple pattern in the grade. Reversion, for example, leaves numbers and arrows alone and changes the sign of planes and volumes.
\emph{First proved in} the algebra core of the Transport Lab, proved there with an exact check; standard, see P. Lounesto, Clifford Algebras and Spinors, 2nd ed., London Mathematical Society Lecture Note Series 286, Cambridge University Press (2001), Ch. 14 (definitions of the Clifford algebra).
\end{theorem}
\noindent\emph{Rests on:} \ref{c:std-algebra.functoriality}, \ref{c:std-algebra.existence-uniqueness}, \ref{c:std-algebra.grade-decomposition}.

The elements that commute with everything form the centre of the algebra.

With $n$ even only the scalars are central. With $n$ odd the pseudoscalar $\omega = e_1 e_2 \cdots e_n$, the product of all the basis vectors, is central too.

Whether the pseudoscalar squares to plus or minus one decides much of the algebra's shape, as the classification below shows.

\begin{lemma}[The centre]\label{c:primer.centre}
The centre of an algebra is the set of its elements that commute with everything in it. For the Clifford algebra of a space with an even number of dimensions, only the ordinary numbers do this. With an odd number of dimensions there is one more such element, the product of all the basis directions, so the centre becomes two-dimensional.
\emph{First proved in} T0, Lem 9.4 (\href{https://doi.org/10.5281/zenodo.22714501}{doi:10.5281/zenodo.22714501}).
\end{lemma}

A vector $u$ with $Q(u) \ne 0$ can be divided by. Sandwiching another vector $x$ between it and its inverse gives \[ x \mapsto -u x u\^{}{-1}, \] the reflection in the mirror perpendicular to $u$.

So the algebra does geometry directly: a mirror is just a vector.

\begin{theorem}[Reflection in a non-null vector]\label{c:std-algebra.reflection}
Multiplying an arrow on both sides by a fixed arrow, with a sign, gives its mirror image in the plane perpendicular to the fixed arrow. This works for any fixed arrow that does not have zero length. Mirrors are the building blocks of all rotations.
\emph{First proved in} the algebra core of the Transport Lab, proved there with an exact check; standard, see C. Doran and A. Lasenby, Geometric Algebra for Physicists, Cambridge University Press (2003), Ch. 4 (foundations of geometric algebra: rotors, versors).
\end{theorem}
\noindent\emph{Rests on:} \ref{c:std-algebra.inverse-of-vector}, \ref{c:std-algebra.vector-squares}, \ref{c:std-algebra.orthogonal-anticommute}.

Two reflections in a row give a rotation, or a boost in spacetime, and the product of the two vectors is a rotor $R$. A rotor acts on vectors by \[ x \mapsto R x \tilde{R} \]

In a nondegenerate space every such rotation arises this way, and exactly two rotors, $R$ and $-R$, give the same one. This doubling is where spinors come from.

\begin{theorem}[Spin double-covers the special orthogonal group]\label{c:std-algebra.spin-double-cover}
Every rotation, and every Lorentz boost, is produced by a rotor, and by exactly two of them, which differ by a sign. This two-to-one relation is the origin of spin one-half: a full turn of space changes the sign of the rotor.
\emph{First proved in} the algebra core of the Transport Lab, proved there with an exact check; standard, see P. Lounesto, Clifford Algebras and Spinors, 2nd ed., London Mathematical Society Lecture Note Series 286, Cambridge University Press (2001), Ch. 17 (spin groups and spinor spaces).
\end{theorem}
\noindent\emph{Rests on:} \ref{c:primer.twisted-adjoint-action}, \ref{c:primer.centre}, \ref{c:std-algebra.cartan-dieudonne}, \ref{c:std-algebra.reflection}, \ref{c:std-algebra.inverse-of-vector}, \ref{c:std-algebra.orthogonal-anticommute}, \ref{c:std-algebra.z2-grading}.

Light-like vectors, nonzero vectors with $Q(u) = 0$, cannot be divided by, but they come in pairs.

A light-like vector $u$ and a partner $w$ span a hyperbolic plane, and the space splits into as many such planes as possible: \[ V = H\_1 \perp \dots \perp H\_m \perp A \perp \operatorname{rad} V, \] with $m = \min(p, q)$ and $A$ free of light-like vectors.

The number of these planes is the Witt index, which later in the lab counts quantum capacity.

\begin{theorem}[Hyperbolic pairs and the Witt decomposition]\label{c:std-algebra.witt-decomposition}
Any space with both space-like and time-like directions can be cut into perpendicular copies of the simplest space-time plane, plus a leftover piece that is purely space-like or purely time-like, plus the directions the length rule cannot see. The number of space-time planes, the smaller of the two counts, is the largest number of independent mutually perpendicular light-like directions the space can hold.
\emph{First proved in} the algebra core of the Transport Lab, proved there with an exact check; standard, see T. Y. Lam, Introduction to Quadratic Forms over Fields, Graduate Studies in Mathematics 67, American Mathematical Society (2005), Ch. I (hyperbolic planes, the Witt decomposition, Witt's theorems).
\end{theorem}
\noindent\emph{Rests on:} \ref{c:std-algebra.sylvester-inertia}, \ref{c:std-algebra.polarisation}, \ref{c:std-algebra.vector-squares}.

Adding one hyperbolic plane doubles the matrix size: \[ \Cl(p{+}1, q{+}1) \cong \Mat\_2(\Cl(p,q)) \]

Adding eight positive or eight negative directions gives sixteen-by-sixteen matrices over the old algebra. So every real geometric algebra is a matrix algebra over $\R$, $\C$ or $\HH$, or a sum of two such, fixed by $n = p + q$ and by $p - q$ modulo $8$.

\begin{theorem}[Periodicity and the table of real Clifford algebras]\label{c:std-algebra.periodicity}
The algebras of all the different spaces fall into a repeating pattern of eight. Each is a table of real, complex or quaternion numbers, sometimes two such tables side by side, and which kind depends only on the difference between the number of space-like and time-like directions, counted round a clock of eight hours.
\emph{First proved in} the algebra core of the Transport Lab, proved there with an exact check; standard, see P. Lounesto, Clifford Algebras and Spinors, 2nd ed., London Mathematical Society Lecture Note Series 286, Cambridge University Press (2001), Ch. 16 (matrix representations and periodicity of 8).
\end{theorem}
\noindent\emph{Rests on:} \ref{c:std-algebra.tensor-rules}, \ref{c:std-algebra.cl11-matrix}, \ref{c:std-algebra.blade-basis}.

Return now to the light-like pairs of a hyperbolic plane.

Chosen as a Witt basis $v_i$, $w_i$, with $B(v_i, w_j) = \tfrac12 \delta_{ij}$, they satisfy \[ v\_i w\_j + w\_j v\_i = \delta\_{ij}, \] and the $v_i$ anticommute among themselves, as do the $w_i$. These are the canonical anticommutation relations, with $v_i$ annihilating and $w_i$ creating a fermion.

One rule about square lengths therefore already contains the algebra of fermions.

\begin{theorem}[Null pairs satisfy the canonical anticommutation relations]\label{c:std-algebra.witt-basis-car}
Inside a Clifford algebra, a direction of space paired with a direction of time produces two zero-length arrows that obey exactly the rules for creating and destroying a fermion. With several such pairs you get several independent fermion modes, without building any Hilbert space first.
\emph{First proved in} the algebra core of the Transport Lab, proved there with an exact check; standard, see C. Doran and A. Lasenby, Geometric Algebra for Physicists, Cambridge University Press (2003), Ch. 9 (multiparticle states).
\end{theorem}
\noindent\emph{Rests on:} \ref{c:std-algebra.vector-squares}.

Let the $w_i$ act by adding a direction and the $v_i$ by removing one. This writes the whole algebra as matrices on all combinations of the $w_i$: \[ \Cl(V, Q) \cong \End(\Lambda W) \]

That space is the fermion Fock space, and it sits inside the algebra as the left ideal generated by the vacuum idempotent $f$, so states and operators live in one algebra.

\begin{theorem}[The Fock representation, and the Fock space as a minimal left ideal]\label{c:std-algebra.fock-representation}
The state space of $k$ fermion modes, with its creation and annihilation operators, already sits inside the Clifford algebra of $k$ space and $k$ time directions. The "empty" state is an element of the algebra, and the other states are obtained by multiplying it by the creation arrows.
\emph{First proved in} the algebra core of the Transport Lab, proved there with an exact check; standard, see C. Chevalley, The Algebraic Theory of Spinors, Columbia University Press (1954), Ch. II and III (spinors, pure spinors).
\end{theorem}
\noindent\emph{Rests on:} \ref{c:std-algebra.witt-basis-car}, \ref{c:std-algebra.existence-uniqueness}, \ref{c:std-algebra.blade-basis}.

\section{Results used from elsewhere}
The proofs above use the following results of the series, stated briefly; each is proved in the paper named.

\begin{fact}[The twisted adjoint action]\label{c:primer.twisted-adjoint-action}
The twisted adjoint action of a versor is an algebra automorphism acting on vectors by hyperplane reflections, giving a two-to-one map of $\Pin(p,q)$ onto $\mathrm{O}(p,q)$.
\emph{Source:} T0, Lem 9.2 (\href{https://doi.org/10.5281/zenodo.22714501}{doi:10.5281/zenodo.22714501}).
\end{fact}

\begin{fact}[The Cartan-Dieudonne theorem]\label{c:std-algebra.cartan-dieudonne}
In a nondegenerate $n$-dimensional real quadratic space every isometry is a product of at most $n$ reflections; this fails for degenerate forms.
\emph{Source:} T. Y. Lam, Introduction to Quadratic Forms over Fields, AMS (2005)., the Cartan-Dieudonne theorem; I. R. Porteous, Clifford Algebras and the Classical Groups, Cambridge University Press (1995)., orthogonal groups generated by reflections.
\end{fact}

\begin{fact}[Cl(1,1) and Cl(2,0) are the real 2x2 matrices]\label{c:std-algebra.cl11-matrix}
$\Cl(H) \cong \Mat_2(\R)$ for a hyperbolic plane, with a hyperbolic pair going to the off-diagonal matrix units; also $\Cl(2,0) \cong \Mat_2(\R)$.
\emph{Source:} the algebra core of the Transport Lab, proved there with an exact check; standard, see P. Lounesto, Clifford Algebras and Spinors, 2nd ed., London Mathematical Society Lecture Note Series 286, Cambridge University Press (2001), Ch. 16 (matrix representations and periodicity of 8).
\end{fact}

\begin{fact}[The Clifford algebra exists and is unique]\label{c:std-algebra.existence-uniqueness}
$\Cl(V, Q) = T(V)/\langle v \otimes v - Q(v) \rangle$ exists and is unique up to unique isomorphism; the wedge-plus-contraction operators on $\Lambda V$ give a concrete representation.
\emph{Source:} the algebra core of the Transport Lab, proved there with an exact check; standard, see C. Chevalley, The Algebraic Theory of Spinors, Columbia University Press (1954), Ch. II (the Clifford algebra of a quadratic space).
\end{fact}

\begin{fact}[Isometries extend to algebra homomorphisms]\label{c:std-algebra.functoriality}
An isometry $f$ extends uniquely to an algebra homomorphism $\Cl(f)$, injective when $f$ is; $\mathrm{O}(V, Q)$ acts by automorphisms.
\emph{Source:} the algebra core of the Transport Lab, proved there with an exact check; standard, see P. Lounesto, Clifford Algebras and Spinors, 2nd ed., London Mathematical Society Lecture Note Series 286, Cambridge University Press (2001), Ch. 14 (definitions of the Clifford algebra).
\end{fact}

\begin{fact}[The grade decomposition]\label{c:std-algebra.grade-decomposition}
In an orthogonal basis $\Cl^{k} = \operatorname{span}\{e_I : |I| = k\}$, so $\Cl = \bigoplus_k \Cl^{k}$ with $\dim \Cl^{k} = \binom{n}{k}$, independently of the basis.
\emph{Source:} the algebra core of the Transport Lab, proved there with an exact check; standard, see C. Doran and A. Lasenby, Geometric Algebra for Physicists, Cambridge University Press (2003), Ch. 4 (foundations of geometric algebra).
\end{fact}

\begin{fact}[Which vectors and versors are invertible]\label{c:std-algebra.inverse-of-vector}
$v$ is invertible iff $Q(v) \ne 0$, with $v^{-1} = v/Q(v)$; a versor $g = u_1 \cdots u_k$ has $g\tilde{g} = \prod Q(u_i)$.
\emph{Source:} the algebra core of the Transport Lab, proved there with an exact check; standard, see C. Doran and A. Lasenby, Geometric Algebra for Physicists, Cambridge University Press (2003), Ch. 4 (foundations of geometric algebra: rotors, versors).
\end{fact}

\begin{fact}[Orthogonal vectors anticommute]\label{c:std-algebra.orthogonal-anticommute}
$uv = -vu$ iff $B(u,v) = 0$; in an orthogonal basis $e_I e_J = \pm \prod_{I \cap J} Q(e_i)\; e_{I \triangle J}$ with the sorting sign.
\emph{Source:} the algebra core of the Transport Lab, proved there with an exact check; standard, see C. Doran and A. Lasenby, Geometric Algebra for Physicists, Cambridge University Press (2003), Ch. 4 (foundations of geometric algebra).
\end{fact}

\begin{fact}[Polarisation]\label{c:std-algebra.polarisation}
The polar form is symmetric with $B(v,v) = Q(v)$, every symmetric bilinear form arises this way, and in a basis $Q(x) = x^{\mathsf T} G x$.
\emph{Source:} the algebra core of the Transport Lab, proved there with an exact check; standard, see T. Y. Lam, Introduction to Quadratic Forms over Fields, Graduate Studies in Mathematics 67, American Mathematical Society (2005), Ch. I and II (quadratic forms; the real case).
\end{fact}

\begin{fact}[Orthogonal bases and the law of inertia]\label{c:std-algebra.sylvester-inertia}
Every real quadratic space has an orthogonal basis squaring to $1, -1, 0$, and the counts $(p, q, r)$ are invariants that classify it up to isometry.
\emph{Source:} T. Y. Lam, Introduction to Quadratic Forms over Fields, AMS (2005)., the diagonalisation of forms and Sylvester's law of inertia (real case).
\end{fact}

\begin{fact}[The tensor product rules]\label{c:std-algebra.tensor-rules}
For $W$ even-dimensional with $\omega_W^{2} = \epsilon$, $\Cl(V \perp W, \epsilon Q \oplus Q_W) \cong \Cl(V) \otimes \Cl(W)$; hence $\Cl(p{+}1,q{+}1) \cong \Mat_2(\Cl(p,q))$, $\Cl(q{+}2,p) \cong \Cl(p,q) \otimes \Cl(2,0)$, $\Cl(q,p{+}2) \cong \Cl(p,q) \otimes \Cl(0,2)$.
\emph{Source:} the algebra core of the Transport Lab, proved there with an exact check; standard, see T. Y. Lam, Introduction to Quadratic Forms over Fields, Graduate Studies in Mathematics 67, American Mathematical Society (2005), Ch. V (Clifford algebras: the tensor product rules).
\end{fact}

\begin{fact}[The even subalgebra and the Z2 grading]\label{c:std-algebra.z2-grading}
$\Cl = \Cl^{+} \oplus \Cl^{-}$ is a $\mathbb{Z}_2$-grading, $\dim \Cl^{+} = 2^{n-1}$, and $\Cl^{j}\Cl^{k}$ has grades $|j-k|, |j-k|+2, \dots, j+k$.
\emph{Source:} the algebra core of the Transport Lab, proved there with an exact check; standard, see P. Lounesto, Clifford Algebras and Spinors, 2nd ed., London Mathematical Society Lecture Note Series 286, Cambridge University Press (2001), Ch. 14 (definitions of the Clifford algebra).
\end{fact}

\section*{Conclusion}
From the single rule $v^{2} = Q(v)$ the walk has reached rotors, the classification of every real geometric algebra, and the fermion Fock space.

The trail rests on no postulate. Each step is proved in the lab and checked by machine, citing the classical law of inertia and the theorem that every isometry is a product of reflections.

\end{document}